\section{任意级数下的抵消条件}
\label{sec:general-order}

\subsection{梯形网络的递推表示}

考虑包含 $N\geq1$ 级的支路，每级均由正值串联电阻 $r_j$ 及其后的正值对地电容 $q_j$ 组成。元件归一化方法和频率变量沿用式~\eqref{eq:3}。图~\ref{fig:circuit} 中的二级电路属于这类网络。取电流从网络流向端接负载为正方向，第 $j$ 级的传输矩阵为
\begin{equation}\label{eq:general-matrix}
M_j(x)=\begin{bmatrix}1+r_jq_jx&r_j\\q_jx&1\end{bmatrix}.
\end{equation}
设 $(A_j,B_j)$ 为 $M_1\cdots M_j$ 的第一行。接入末级后，有
\begin{equation}\label{eq:general-forward}
B_j=r_jA_{j-1}+B_{j-1},\qquad
A_j=A_{j-1}+xq_jB_j,
\quad A_0=1,\ B_0=0.
\end{equation}
由式~\eqref{eq:general-forward} 归纳可得 $\deg A_N=N$、$\deg B_N=N-1$ 且 $A_N(0)=1$。两个多项式的最高次项系数均为正。各级矩阵的行列式均为 1，其乘积的行列式也为 1，因此第一行的两个多项式不存在非常数公因子。

由端口方程可得开路传递函数和输入源短路时的输出阻抗。在两条支路分别接入相同电阻负载后，即得到下述带载传递函数；该端接分解与经典综合中的结构一致\citep{guillemin1949rcsynthesis}：
\begin{equation}\label{eq:general-transfer}
H_0=\frac1{A_N},\qquad
\frac{Z_0}{R_*}=\frac{B_N}{A_N},\qquad
H_g=\frac1{A_N+gB_N}.
\end{equation}
当 $N=2$ 时，$A_2=1+ax+bx^2$、$B_2=c+dx$，即得到式~\eqref{eq:5}。保留 $B_N$ 可描述开路电压传递函数未包含的端接信息。

\subsection{经典梯形提取与系数坐标}

输入源短路后，$B_N/A_N$ 为从输出端看入的无量纲驱动点阻抗，可写成梯形连分式\citep{guillemin1949rcsynthesis,lee1978chain}
\[
z_j=\left[q_jx+\left(r_j+z_{j-1}\right)^{-1}\right]^{-1},
\qquad z_0=0,\quad z_N=B_N/A_N.
\]
下面的逐级恢复将已有连分式梯形综合\citep{guillemin1949rcsynthesis,lee1978chain} 特化为单电阻/单电容臂拓扑的 $(A_j,B_j)$ 坐标；一般臂梯形网络的递推传递计算也属经典方法\citep{fialkow1951rcladder}。本文利用该提取结果，计算指定端接约束下的可实现元件族维数。

末级元件值可由最高次项系数恢复。以 $\operatorname{lc}(p)$ 表示多项式 $p$ 的最高次项系数，定义
\begin{equation}\label{eq:general-inverse}
\begin{aligned}
q_j&=\frac{\operatorname{lc}(A_j)}{\operatorname{lc}(B_j)},&
A_{j-1}&=A_j-xq_jB_j,\\
r_j&=\frac{\operatorname{lc}(B_j)}{\operatorname{lc}(A_{j-1})},&
B_{j-1}&=B_j-r_jA_{j-1}.
\end{aligned}
\end{equation}
对于全正值元件梯形网络，上述运算恰好逆转式~\eqref{eq:general-forward} 的递推。每一步中，$A_{j-1}$ 的次数为 $j-1$；当 $j>1$ 时，$B_{j-1}$ 的次数为 $j-2$。最终得到 $(A_0,B_0)=(1,0)$。各比值均为正值，且唯一确定。

反之，从满足上述次数要求且 $A_N(0)=1$ 的实系数多项式出发，若每次相减均得到要求的次数、恢复出的每个元件值均为正，且最终得到 $(1,0)$，则逆序执行上述运算，即可由全正值元件梯形网络重构原多项式。因此，这些条件给出了该拓扑可实现性的有限步充要检验。正值或次数检查中任一项不满足，即可排除该候选系数对。

从 $2N$ 个正值元件到 $A_N$ 和 $B_N$ 的 $2N$ 个系数的正向映射为多项式映射。式~\eqref{eq:general-inverse} 给出的逆映射为有理映射，在每个正值元件实现处，其分母均非零。在足够小的邻域内，元件值仍严格为正，最高次项系数仍非零。两个映射在该邻域内光滑且互逆，因此这些系数可作为元件族的局部坐标。

\subsection{匹配自由度的数目}

固定一条参考支路。开路匹配固定了 $A_N$ 的 $N$ 个非常数项系数，在正值元件可实现域内保留 $N$ 个局部自由度。同一传递函数具有多个正值实现属经典结论\citep{fialkow1951rcladder,fialkow1952rcnetworks}；此处计数的是本拓扑内与固定参考支路匹配的元件族。统一阻抗缩放使 $B_N$ 乘以同一个标量因子，同时保持 $A_N$ 不变。当 $N\geq2$ 时，该缩放仅沿更大匹配族中的一个方向变化。

在指定的 $g>0$ 下，原始响应匹配 $H_{g,P}=H_{g,N}$ 固定 $D_g=A_N+gB_N$ 的全部系数。已有综合理论以限制可实现增益范围的方式处理端接\citep{fialkow1951rcladder,fialkow1952rcnetworks}；以下计数针对的则是与给定参考支路匹配的元件族。常数项通过 $D_g(0)=1+gB_N(0)$ 固定 $B_N(0)$，其余 $N$ 个系数固定 $A_N$ 与 $B_N$ 系数的组合。这 $N+1$ 个独立线性约束保留 $N-1$ 个局部自由度，$N=1$ 时为零。

直流校准后比较的对象为 $G_g=D_g(0)/D_g$。固定归一化分母产生 $N$ 个独立约束，保留 $N$ 个局部自由度。在系数坐标中，可以通过改变 $A_N$，独立改变归一化分母的最高次项及其余各项系数，故约束映射的秩为 $N$。此处使用的增益自由度在指定增益综合中已有对应结论：固定极点后，增益常数可在可实现区间内变化\citep{fialkow1951rcladder,paige1960maximum,wang2017rcladder}。在本文的匹配问题中，该自由度表现为直流校准后配对支路的歧义，并影响未经校准的差分输出。对于 $N=1$，正值元件族
\[
r>0,\qquad q=\frac{r+10}{11r}
\]
在 $g=0.1$ 时与 $(r,q)=(1,1)$ 具有相同归一化响应，但其原始直流增益随 $r$ 改变。

对于 $N=2$，支路 P = $(1,1,1,1)$ 与正值支路
\[
\left(2,2,\frac{217+\sqrt{6769}}{288},\frac{217-\sqrt{6769}}{480}\right)
\]
在 $g=0.1$ 时具有相同归一化响应。两者原始分母相差常数倍 $7/6$。写成 $D_N=(7/6)D_P$ 可得 $H_N=(6/7)H_P$，故原始共模转差模增益为 $G_{DC}=H_P-H_N=H_P/7$，差模增益为 $G_{DD}=(H_P+H_N)/2=(13/14)H_P$。两者之比与频率无关，原始 CMRR 在所有频率均为 $20\log_{10}|G_{DD}/G_{DC}|=20\log_{10}(13/2)=16.26$~dB。校准因此允许额外的匹配方向，该方向并不消除原始差分输出：读出所用的归一化响应一致，而原始共模转差模转换仍然存在。

\begin{table}[htbp]
\centering
\caption{固定正值 $N$ 级参考支路时，电阻端接下的局部匹配维数。原始匹配包含绝对增益，归一化匹配比较 $H_g/H_g(0)$。}
\label{tab:matching-conventions}
\begin{tabular}{lll}
\toprule
端接信息 & 原始响应 & 归一化响应 \\
\midrule
开路 & $N$ & $N$ \\
一个指定 $g>0$ & $N-1$ & $N$ \\
两个不同指定 $g_1,g_2>0$ & $0$ & $0$ \\
\bottomrule
\end{tabular}
\end{table}

对于两个不同的正值电导，令 $U_i=1/H_{g_i}$。两个完整传递函数可确定
\begin{equation}\label{eq:general-two-load}
B_N=\frac{U_2-U_1}{g_2-g_1},\qquad
A_N=\frac{g_2U_1-g_1U_2}{g_2-g_1}.
\end{equation}
因此，若两条支路在两种负载下均匹配，则其 $A_N$ 和 $B_N$ 分别相同。由经典梯形提取的逐步唯一性可知，对应元件值必须相同\citep{guillemin1949rcsynthesis}。该论证适用于上述梯形网络类别中的任意有限 $N$，所用信息为两种负载下包含绝对增益的完整复传递函数。

两个归一化电阻端接响应也可精确确定元件。对候选支路，写成 $D_i(g_k)=c_kD_{\rm ref}(g_k)$。$x^N$ 的系数不随 $g_k$ 改变且非零，因此 $c_1=c_2=c$。两负载方程相减得 $B_i=cB_{\rm ref}$，进而得 $A_i=cA_{\rm ref}$。由 $A_i(0)=A_{\rm ref}(0)=1$ 可知 $c=1$。该证明使用精确完整响应，尚未规定含噪数据的稳定估计方法。

\subsection{负载类型对匹配族的影响}

已知的有理函数负载可用无量纲导纳 $g(x)=P(x)/Q(x)$ 表示，其中 $P$ 和 $Q$ 为非零且互素的实系数多项式。假设 $QA_N+PB_N$ 不恒为零，使带载响应有定义。两条支路的传递函数相等可得
\begin{equation}\label{eq:rational-kernel}
H_g=\frac{Q}{QA_N+PB_N},\qquad
Q\Delta A+P\Delta B=0.
\end{equation}
其中 $\Delta$ 表示支路 P 减去支路 N。由互素性可知，存在多项式 $S$，使 $\Delta A=PS$ 且 $\Delta B=-QS$。次数上限和 $\Delta A(0)=0$ 对 $S$ 的约束为
\begin{equation}\label{eq:rational-dimension}
\begin{aligned}
m&=\min\{N-\deg P,\ N-1-\deg Q\},\\
d_g&=\max\{0,m+1-\epsilon\},\qquad
\epsilon=\begin{cases}1,&P(0)\ne0,\\0,&P(0)=0.\end{cases}
\end{aligned}
\end{equation}
当 $m<0$ 时，仅允许 $S=0$。否则，$S$ 有 $m+1$ 个系数；若 $P(0)$ 非零，则其常数项系数必须为零。该线性族与正值元件可实现开域相交后，局部维数仍为 $d_g$。

表~\ref{tab:load-dimensions} 给出将式~\eqref{eq:rational-dimension} 用于常见负载的结果。开路时 $P$ 恒为零，单独列出。纯电容负载改变端接电容，保留 $N$ 个自由度；电阻负载还约束零频增益，将自由度减少至 $N-1$。
\begin{table}[htbp]
\centering
\caption{在一种指定负载下、固定参考支路时的原始响应匹配族局部维数，包含绝对增益约束。表中各负载参数均为正值。直流归一化匹配另见表~\ref{tab:matching-conventions}。}
\label{tab:load-dimensions}
\begin{tabular}{lll}
\toprule
负载 & 无量纲导纳 & 维数 \\
\midrule
开路 & $0$ & $N$ \\
电阻 & $g_0$ & $N-1$ \\
电容 & $\alpha x$ & $N$ \\
电感 & $\beta/x$ & $\max(N-2,0)$ \\
串联 RC & $\alpha x/(1+\gamma x)$ & $N-1$ \\
并联 RC & $g_0+\alpha x$ & $N-1$ \\
\bottomrule
\end{tabular}
\end{table}

在原始响应匹配口径下，两个不同的已知有理导纳也可确定元件组合。令 $g_k=P_k/Q_k$，分别将 $Q_k\Delta A+P_k\Delta B=0$ 乘以另一负载的分母后相减，得到 $(P_1Q_2-P_2Q_1)\Delta B=0$。负载函数不恒等意味着 $P_1Q_2-P_2Q_1\not\equiv0$，因此 $\Delta B=0$，进而 $\Delta A=0$。已知非零多项式 $Q_k$ 可在有理函数恒等式中约去，孤立零点处由延拓处理。表~\ref{tab:load-dimensions} 中的负载均为时不变有理导纳；周期开关输入需要第~\ref{sec:dynamic} 节另行引入的采样模型。

\subsection{负载变化与恢复对测量误差的敏感性}

对于独立支路，令 $D_i(g)=A_i+gB_i$。共模转差模增益的精确表达式为
\begin{equation}\label{eq:general-leakage}
G_{DC}=-\frac{\Delta A+g\Delta B}{D_P(g)D_N(g)}.
\end{equation}
若两条支路在 $g_0>0$ 下匹配，则 $\Delta A=-g_0\Delta B$ 且 $\Delta B(0)=0$。当两侧同时换用另一种相同电阻负载时，
\begin{equation}\label{eq:general-load-shift}
G_{DC}(g)=-\frac{(g-g_0)\Delta B}{D_P(g)D_N(g)}.
\end{equation}
零频泄漏仍为零，其余随频率变化的失配由多项式 $\Delta B$ 决定。二阶网络的负载变化公式对应 $\Delta B=x\Delta d$。

精确可辨识性与恢复对测量误差的敏感性需要分别检查。式~\eqref{eq:general-two-load} 以负载电导之差为分母，响应的小扰动则满足 $\delta(1/H)\simeq-\delta H/H^2$。两种负载过于接近或响应衰减过强，均会放大测量误差。因此，实验中的元件恢复除代数唯一性结论外，还需要明确频带和不确定性模型。
对无量纲系数映射 $(r_1,\ldots,r_N,q_1,\ldots,q_N)\mapsto(A_{N,1},\ldots,A_{N,N},B_{N,0},\ldots,B_{N,N-1})$，单位元件点未缩放Jacobian的2范数条件数在 $N=2,4,6,8,10,12$ 时分别为 $15.6$、$1.08\times10^3$、$9.30\times10^4$、$8.67\times10^6$、$8.40\times10^8$、$8.35\times10^{10}$。这些数值针对所列系数归一化，随坐标缩放改变，描述系数恢复的局部误差放大，并非特定响应拟合估计器的性能界。
\FloatBarrier
